\documentclass[11pt,a4paper]{article}
% LTeX: language=fr
% Character encoding and font configuration for standard pdflatex
\usepackage[utf8]{inputenc}
\usepackage[T1]{fontenc}
\usepackage{mathptmx} % Times New Roman clone for a professional academic look
\usepackage{graphicx}
\usepackage[left=2.5cm,right=2.5cm,top=2cm,bottom=2cm]{geometry}
\usepackage{parskip}   % Automatically adds space between paragraphs and removes indentation
\usepackage[colorlinks=true, urlcolor=black, linkcolor=black]{hyperref}
\usepackage{circuitikz}
\usepackage{float}

\providecommand{\SchemaParalleles}{
\begin{circuitikz}
\tikzstyle{every node}=[font=\fontsize{10.2pt}{13.3pt}\selectfont]
\draw [dash dot dot] (17.75,14) -- (22.125,14);
\draw  (19.125,10.75) rectangle (20.75,10.25);
\draw  (22.25,14.25) rectangle (23.875,13.75);
\node [font=\fontsize{10.2pt}{13.3pt}\selectfont, inner xsep=0.080cm, inner ysep=0.085cm, rounded corners=0.020cm] at (24.25,14) {Fixe};
\node [font=\fontsize{14.8pt}{19.2pt}\selectfont, inner xsep=0.080cm, inner ysep=0.085cm, rounded corners=0.020cm] at (23.125,14.625) {$UWB_1$};
\node [font=\fontsize{14.8pt}{19.2pt}\selectfont, inner xsep=0.080cm, inner ysep=0.085cm, rounded corners=0.020cm] at (20,9.875) {$UWB_2$};
\draw [short] (23.25,13.25) -- (23.25,13.25);
\draw [short] (22.75,13.75) -- (22.75,13.75);
\draw [dotted] (23.125,13.625) -- (23.125,10);
\draw [-{Stealth[open,scale=1.5]}, dash dot dot] (20.875,10.5) -- (27,10.5);
\draw [short] (23,13.125) -- (23,13.125);
\draw [short] (23,12.75) -- (23,12.75);
\draw [densely dotted] (22.875,11.75) -- (22.875,11.75);
\draw [densely dotted] (23,12.125) -- (23,12.125);
\draw [densely dotted] (23,12.125) -- (23,12.125);
\draw [{Stealth[scale=1.5]}-{Stealth[scale=1.5]}, ] (20,10.875) -- (23,13.625);
\draw [short] (23,12.25) arc[start angle=-74.74, end angle=-164.74, radius=0.7126];
\node [font=\fontsize{10.2pt}{13.3pt}\selectfont, inner xsep=0.080cm, inner ysep=0.085cm, rounded corners=0.000cm] at (22.375,12.125) {$\theta$};
\draw [dash dot dot] (24.625,14) -- (27,14);
\draw [{Stealth[open,scale=1.5]}-, dash dot dot] (17.625,10.5) -- (19,10.5);
\node [font=\fontsize{10.2pt}{13.3pt}\selectfont, inner xsep=0.080cm, inner ysep=0.085cm, rounded corners=0.000cm] at (23.375,10.25) {y};
\draw [short] (21.375,11.125) -- (21.375,11.125);
\draw [short] (20.75,11.625) -- (20.75,11.625);
\draw [line width=1.2pt, short] (19.625,10.75) -- (20.25,10.75);
\draw [line width=1.2pt, short] (22.75,13.75) -- (23.375,13.75);
\end{circuitikz}
}

\providecommand{\Schemarayon}{
\begin{circuitikz}
\draw [dash dot dot] (25.375,7.875) -- (29.75,7.875);
\draw [ rotate around={-45:(27.3125, 4.5)}] (26.5,4.75) rectangle (28.125,4.25);
\draw  (29.875,8.125) rectangle (31.5,7.625);
\node [font=\fontsize{10.2pt}{13.3pt}\selectfont, inner xsep=0.080cm, inner ysep=0.085cm, rounded corners=0.020cm] at (31.875,7.875) {Fixe};
\node [font=\fontsize{14.8pt}{19.2pt}\selectfont, inner xsep=0.080cm, inner ysep=0.085cm, rounded corners=0.020cm] at (30.75,8.5) {$UWB_1$};
\node [font=\fontsize{14.8pt}{19.2pt}\selectfont, rotate around={-45:(0,0)}, inner xsep=0.080cm, inner ysep=0.085cm, rounded corners=0.020cm] at (26.875,4) {$UWB_2$};
\draw [short] (30.875,7.125) -- (30.875,7.125);
\draw [short] (30.375,7.625) -- (30.375,7.625);
\draw [dotted] (30.75,7.5) -- (30.75,3.875);
\draw [short] (30.625,7) -- (30.625,7);
\draw [short] (30.625,6.625) -- (30.625,6.625);
\draw [densely dotted] (30.5,5.625) -- (30.5,5.625);
\draw [densely dotted] (30.625,6) -- (30.625,6);
\draw [densely dotted] (30.625,6) -- (30.625,6);
\draw [{Stealth[scale=1.5]}-{Stealth[scale=1.5]}, ] (27.625,4.75) -- (30.625,7.5);
\draw [short] (30.625,5.75) arc[start angle=-74.74, end angle=-164.74, radius=0.7126];
\node [font=\fontsize{10.2pt}{13.3pt}\selectfont, inner xsep=0.080cm, inner ysep=0.085cm, rounded corners=0.000cm] at (30,6) {$\theta$};
\draw [dash dot dot] (32.25,7.875) -- (34.625,7.875);
\node [font=\fontsize{10.2pt}{13.3pt}\selectfont, inner xsep=0.080cm, inner ysep=0.085cm, rounded corners=0.000cm] at (31,4.125) {y};
\draw [short] (29,5) -- (29,5);
\draw [short] (28.375,5.5) -- (28.375,5.5);
\draw [line width=1.2pt, short] (27.25,4.875) -- (27.75,4.375);
\draw [line width=1.2pt, short] (30.375,7.625) -- (31,7.625);
\node [font=\fontsize{14.8pt}{19.2pt}\selectfont, rotate around={41:(0,0)}, inner xsep=0.080cm, inner ysep=0.085cm, rounded corners=0.010cm] at (29,6.375) {$d$};
\draw [dash dot dot] (30.625,2.75) arc[start angle=-95.88, end angle=-130.52, radius=4.7969];
\draw [dash dot dot] (26.625,5.25) arc[start angle=-138.76, end angle=-173.32, radius=4.1447];
\end{circuitikz}
}

\providecommand{\FtFhorizontale}{
\begin{circuitikz}
\tikzstyle{every node}=[font=\fontsize{10.2pt}{13.3pt}\selectfont]
    \draw  (32.5,13.5) rectangle (34.75,13.25);
    \draw  (34.125,13.625) rectangle (34.75,13.5);
    \draw  (32.875,13.625) rectangle (33.75,13.5);
    \draw  (38.375,13.5) rectangle (40.625,13.25);
    \draw  (39.5,13.625) rectangle (40.375,13.5);
    \draw  (38.375,13.625) rectangle (39,13.5);
    \draw [{Stealth[scale=1.5]}-{Stealth[scale=1.5]}, ] (35,13.5) -- (38.125,13.5);
    \node [font=\fontsize{18.2pt}{23.7pt}\selectfont, inner xsep=0.080cm, inner ysep=0.085cm, rounded corners=0.020cm] at (36.5,13.875) {d};
    \node [font=\fontsize{14.8pt}{19.2pt}\selectfont, inner xsep=0.080cm, inner ysep=0.085cm, rounded corners=0.020cm] at (33.75,12.875) {$UWB_1$};
    \node [font=\fontsize{14.8pt}{19.2pt}\selectfont, inner xsep=0.080cm, inner ysep=0.085cm, rounded corners=0.020cm] at (39.625,12.875) {$UWB_2$};
\end{circuitikz}
}


\providecommand{\FtFVertical}{
\begin{circuitikz}
\tikzstyle{every node}=[font=\fontsize{10.2pt}{13.3pt}\selectfont]
    \node [font=\fontsize{14.8pt}{19.2pt}\selectfont, inner xsep=0.080cm, inner ysep=0.085cm, rounded corners=0.020cm] at (33.75,12.875) {$UWB_1$};
    \node [font=\fontsize{14.8pt}{19.2pt}\selectfont, inner xsep=0.080cm, inner ysep=0.085cm, rounded corners=0.020cm] at (39.625,12.875) {$UWB_2$};
    \draw [ rotate around={-90:(36.25, 11.25)}] (35.125,11.375) rectangle (37.375,11.125);
    \draw [ rotate around={-90:(36.4375, 10.9375)}] (36,11) rectangle (36.875,10.875);
    \draw [ rotate around={-90:(36.4375, 12.0625)}] (36.125,12.125) rectangle (36.75,12);
    \draw [ rotate around={-90:(40.375, 11.25)}] (39.25,11.375) rectangle (41.5,11.125);
    \draw [ rotate around={-90:(40.1875, 10.9375)}] (39.75,11) rectangle (40.625,10.875);
    \draw [ rotate around={-90:(40.1875, 12.0625)}] (39.875,12.125) rectangle (40.5,12);
    \draw [{Stealth[scale=1.5]}-{Stealth[scale=1.5]}, ] (36.75,12) -- (39.875,12);
    \node [font=\fontsize{18.2pt}{23.7pt}\selectfont, inner xsep=0.080cm, inner ysep=0.085cm, rounded corners=0.020cm] at (38.25,12.375) {d};
    \node [font=\fontsize{14.8pt}{19.2pt}\selectfont, inner xsep=0.080cm, inner ysep=0.085cm, rounded corners=0.020cm] at (41.375,11.25) {$UWB_2$};
    \node [font=\fontsize{14.8pt}{19.2pt}\selectfont, inner xsep=0.080cm, inner ysep=0.085cm, rounded corners=0.020cm] at (35.375,11.25) {$UWB_1$};
\end{circuitikz}
}

\providecommand{\FtB}{
\begin{circuitikz}
\tikzstyle{every node}=[font=\fontsize{10.2pt}{13.3pt}\selectfont]
    \draw [ rotate around={-90:(36.25, 11.25)}] (35.125,11.375) rectangle (37.375,11.125);
    \draw [ rotate around={-90:(36.4375, 10.9375)}] (36,11) rectangle (36.875,10.875);
    \draw [ rotate around={-90:(36.4375, 12.0625)}] (36.125,12.125) rectangle (36.75,12);
    \draw [{Stealth[scale=1.5]}-{Stealth[scale=1.5]}, ] (36.75,12) -- (39.875,12);
    \node [font=\fontsize{18.2pt}{23.7pt}\selectfont, inner xsep=0.080cm, inner ysep=0.085cm, rounded corners=0.020cm] at (38.25,12.375) {d};
    \node [font=\fontsize{14.8pt}{19.2pt}\selectfont, inner xsep=0.080cm, inner ysep=0.085cm, rounded corners=0.020cm] at (41.375,11.25) {$UWB_2$};
    \node [font=\fontsize{14.8pt}{19.2pt}\selectfont, inner xsep=0.080cm, inner ysep=0.085cm, rounded corners=0.020cm] at (35.375,11.25) {$UWB_1$};
    \draw [ rotate around={-90:(40.25, 11.25)}] (39.125,11.375) rectangle (41.375,11.125);
    \draw [ rotate around={-90:(40.4375, 10.9375)}] (40,11) rectangle (40.875,10.875);
    \draw [ rotate around={-90:(40.4375, 12.0625)}] (40.125,12.125) rectangle (40.75,12);
\end{circuitikz}
}

\begin{document}

\begin{minipage}[c]{0.30\textwidth}
    \includegraphics[width=\linewidth,keepaspectratio]{Image/Logo_INSA.png}
\end{minipage}
\hfill
\begin{minipage}[c]{0.80\textwidth}
    {\LARGE \textbf{Projet S-IoT : Essaim de drone}}\\[2pt]
    {\large Protocole de calibration }
\end{minipage}

\vspace{0.4cm}
\hrule height 1pt
\vspace{0.8cm}

\begin{center}
    {\huge \textbf{Protocole de Test}} \\[5pt]
    {\large\textit{du 09/10/2026}} \\[5pt]
\end{center}

\tableofcontents
\newpage
Le but de ce document est de définir les tests qui ont été fais sur les modules UWB DW1000. Et ensuite, d'exposer les résultats et les conclusions

Les tests sont divisés en deux parties : Sans drone puis avec Drone.

La quasi-totalité des tests ont été réalisés en salle topographie. 
\section{Sans drone}
\subsection{Orientation du module}
Dans cette partie, 
\subsubsection{Test à plat}
Le test à plats a été réaliser un peu différement des autres car effectuer en Plateforme Génie électrique. 
Voici les résultats qui ont été trouvés : 

\begin{table}[H]
    \centering
    \begin{tabular}{|l|l|l|}
    \hline
        Distance réelle (cm)  & Distance mesurée & Erreur \\ \hline
        69,4 & 61 & -8,4 \\ \hline
        70 & 62,2 & -7,8 \\ \hline
        75 & 69,6 & -5,4 \\ \hline
        80 & 73,9 & -6,1 \\ \hline
        82,6 & 76,1 & -6,5 \\ \hline
        85 & 79,9 & -5,1 \\ \hline
        90 & 86,6 & -3,4 \\ \hline
        91 & 87 & -4 \\ \hline
        95 & 92,2 & -2,8 \\ \hline
        100 & 96,5 & -3,5 \\ \hline
        100 & 96,8 & -3,2 \\ \hline
        105 & 102,6 & -2,4 \\ \hline
        107,3 & 106,4 & -0,9 \\ \hline
        110 & 108,9 & -1,1 \\ \hline
        115 & 113,4 & -1,6 \\ \hline
        120,5 & 120,9 & 0,4 \\ \hline
        125 & 126,4 & 1,4 \\ \hline
        130 & 129,6 & -0,4 \\ \hline
        132 & 133,7 & 1,7 \\ \hline
        135 & 133,6 & -1,4 \\ \hline
        140 & 138,6 & -1,4 \\ \hline
        145 & 144,9 & -0,1 \\ \hline
        147,4 & 150 & 2,6 \\ \hline
        150 & 151,7 & 1,7 \\ \hline
        155 & 160,3 & 5,3 \\ \hline
        155,5 & 161 & 5,5 \\ \hline
        160 & 165,7 & 5,7 \\ \hline
        164,7 & 171,8 & 7,1 \\ \hline
        165 & 172,2 & 7,2 \\ \hline
        167 & 172,8 & 5,8 \\ \hline
        170 & 177,2 & 7,2 \\ \hline
        175 & 180,3 & 5,3 \\ \hline
        176 & 183,8 & 7,8 \\ \hline
        180 & 184,3 & 4,3 \\ \hline
        185 & 188,3 & 3,3 \\ \hline
        189,6 & 197,9 & 8,3 \\ \hline
        190 & 191,9 & 1,9 \\ \hline
    \end{tabular}
    \caption{Mesure face à face de façon horizontale}
\end{table}

\subsubsection{Test en vertical}

\begin{table}[H]
    \centering
    \begin{tabular}{|l|l|l|}
    \hline
        Distance réelle (cm) & Distance mesurée (cm) & Erreur  \\ \hline
        70 & 67 & 3 \\ \hline
        80 & 78 & 2 \\ \hline
        90 & 89,2 & 0,8 \\ \hline
        100 & 100 & 0 \\ \hline
        110 & 110,6 & -0,6 \\ \hline
        120 & 120 & 0 \\ \hline
        130 & 132,4 & -2,4 \\ \hline
        140 & 143,6 & -3,6 \\ \hline
        150 & 154,2 & -4,2 \\ \hline
        160 & 165,7 & -5,7 \\ \hline
        170 & 176,4 & -6,4 \\ \hline
    \end{tabular}
    \caption{Mesures pour module face à face avec une référence posé à 120cm}
\end{table}

\begin{table}[H]
    \centering
    \begin{tabular}{|l|l|l|}
    \hline
        Distance réelle (cm) & Distance mesurée (cm) & Erreur  \\ \hline
        70 & 64 & 6 \\ \hline
        80 & 76,2 & 3,8 \\ \hline
        90 & 85,6 & 4,4 \\ \hline
        100 & 97,4 & 2,6 \\ \hline
        110 & 108,2 & 1,8 \\ \hline
        120 & 119,6 & 0,4 \\ \hline
        130 & 130 & 0 \\ \hline
        140 & 140 & 0 \\ \hline
        150 & 152,7 & -2,7 \\ \hline
        160 & 162 & -2 \\ \hline
        170 & 172,2 & -2,2 \\ \hline
    \end{tabular}
    \caption{Mesures pour module face à face avec une référence posé à 130cm}
\end{table}
\subsubsection{Conclusion sur l'orientation}
D'après les essais précédents et les résultats, il est visible que l'orientation verticale semble avoir de meilleure performance de façon globale que le module à plat. 
Dans la suite des essais, nous serons donc toujours avec un module vertical.
\subsection{Influence de l'angle}
Nous avons cherché à définir un semblant de diagramme de rayonnement, qui nous permettrait de déterminer la position optimale des drones. Pour cela, nous avons réalisé deux tests d'angles.

\subsubsection{Module parallèle}
\begin{figure}[H]
  \centering
  \SchemaParalleles
  \caption{Schéma pour les essais des modules en parallèle}
\end{figure}

\begin{table}[H]
    \centering
    \begin{tabular}{|l|l|l|l|l|}
    \hline
        Décalage (cm) & Angle(°) & Distance réelle (cm) & Distance mesurée (cm) & Erreur  \\ \hline
        0 & 0 & 120 & 117 & 3 \\ \hline
        10 & 4,78 & 121 & 118,4 & 2,6 \\ \hline
        20 & 9,59 & 121,5 & 118,9 & 2,6 \\ \hline
        30 & 14,48 & 122,5 & 119,1 & 3,4 \\ \hline
        40 & 19,47 & 125,8 & 124,8 & 1 \\ \hline
        50 & 24,62 & 128,5 & 128,7 & -0,2 \\ \hline
        60 & 30,00 & 132,5 & 134,4 & -1,9 \\ \hline
        70 & 35,69 & 137 & 142,6 & -5,6 \\ \hline
    \end{tabular}
    \caption{Mesure parallèle}
\end{table}

\subsubsection{Sur un rayon}
\begin{figure}[H]
  \centering
  \Schemarayon
  \caption{Schéma pour les essais sur un rayon}
\end{figure}

\begin{table}[H]
    \centering
    \begin{tabular}{|l|l|l|}
    \hline
        Angle(°) & Distance réelle (cm) & Distance mesurée (cm) \\ \hline
        10 & 120 & 116,7 \\ \hline
        20 & 120 & 118,4 \\ \hline
        30 & 120 & 119,5 \\ \hline
        40 & 120 & 122,8 \\ \hline
        50 & 120 & 125,9 \\ \hline
        60 & 120 & 124,2 \\ \hline
        70 & 120 & 123 \\ \hline
        80 & 120 & 124 \\ \hline
        90 & 120 & 130 \\ \hline
    \end{tabular}
    \caption{Mesure rayon}
\end{table}

\subsection{Avant/arrière}

\begin{table}[H]
    \centering
    \begin{tabular}{|l|l|l|}
    \hline
        Distance réelle (cm) & Distance mesurée (cm) & Erreur  \\ \hline
        70 & 97 & -27 \\ \hline
        80 & 109,8 & -29,8 \\ \hline
        90 & 122,8 & -32,8 \\ \hline
        100 & 133,3 & -33,3 \\ \hline
        110 & 143,5 & -33,5 \\ \hline
        120 & 152 & -32 \\ \hline
        130 & 161,6 & -31,6 \\ \hline
        140 & 177,2 & -37,2 \\ \hline
        150 & 188 & -38 \\ \hline
        160 & 192,2 & -32,2 \\ \hline
        170 & 208,8 & -38,8 \\ \hline
    \end{tabular}
    \caption{Mesure face dos }
\end{table}

\section{Sur drone}

\subsection{Orientation du module}
Aucun test à plat n'a été réalisé car prouvez moins bien lors des essais sans drone, par efficacité de temps, le module posé verticalement a été conservé. 
\subsubsection{Test en vertical}

\begin{table}[H]
    \centering
    \begin{tabular}{|l|l|l|}
    \hline
        Distance réelle (cm) & Distance mesurée (cm) & Erreur  \\ \hline
        100 & 92,6 & 7,4 \\ \hline
        110 & 104,3 & 5,7 \\ \hline
        120 & 115 & 5 \\ \hline
        130 & 124,1 & 5,9 \\ \hline
        140 & 134,6 & 5,4 \\ \hline
    \end{tabular}
    \caption{Mesure sur drone face à face  }
\end{table}
\subsection{Influence de l'angle}
Les essais sont faits à 1m20 de distance. 
\subsubsection{Module parallèle}
\providecommand{\SchemaParalleles}{
\begin{circuitikz}
\tikzstyle{every node}=[font=\fontsize{10.2pt}{13.3pt}\selectfont]
\draw [dash dot dot] (17.75,14) -- (22.125,14);
\draw  (19.125,10.75) rectangle (20.75,10.25);
\draw  (22.25,14.25) rectangle (23.875,13.75);
\node [font=\fontsize{10.2pt}{13.3pt}\selectfont, inner xsep=0.080cm, inner ysep=0.085cm, rounded corners=0.020cm] at (24.25,14) {Fixe};
\node [font=\fontsize{14.8pt}{19.2pt}\selectfont, inner xsep=0.080cm, inner ysep=0.085cm, rounded corners=0.020cm] at (23.125,14.625) {$UWB_1$};
\node [font=\fontsize{14.8pt}{19.2pt}\selectfont, inner xsep=0.080cm, inner ysep=0.085cm, rounded corners=0.020cm] at (20,9.875) {$UWB_2$};
\draw [short] (23.25,13.25) -- (23.25,13.25);
\draw [short] (22.75,13.75) -- (22.75,13.75);
\draw [dotted] (23.125,13.625) -- (23.125,10);
\draw [-{Stealth[open,scale=1.5]}, dash dot dot] (20.875,10.5) -- (27,10.5);
\draw [short] (23,13.125) -- (23,13.125);
\draw [short] (23,12.75) -- (23,12.75);
\draw [densely dotted] (22.875,11.75) -- (22.875,11.75);
\draw [densely dotted] (23,12.125) -- (23,12.125);
\draw [densely dotted] (23,12.125) -- (23,12.125);
\draw [{Stealth[scale=1.5]}-{Stealth[scale=1.5]}, ] (20,10.875) -- (23,13.625);
\draw [short] (23,12.25) arc[start angle=-74.74, end angle=-164.74, radius=0.7126];
\node [font=\fontsize{10.2pt}{13.3pt}\selectfont, inner xsep=0.080cm, inner ysep=0.085cm, rounded corners=0.000cm] at (22.375,12.125) {$\theta$};
\draw [dash dot dot] (24.625,14) -- (27,14);
\draw [{Stealth[open,scale=1.5]}-, dash dot dot] (17.625,10.5) -- (19,10.5);
\node [font=\fontsize{10.2pt}{13.3pt}\selectfont, inner xsep=0.080cm, inner ysep=0.085cm, rounded corners=0.000cm] at (23.375,10.25) {y};
\draw [short] (21.375,11.125) -- (21.375,11.125);
\draw [short] (20.75,11.625) -- (20.75,11.625);
\draw [line width=1.2pt, short] (19.625,10.75) -- (20.25,10.75);
\draw [line width=1.2pt, short] (22.75,13.75) -- (23.375,13.75);
\end{circuitikz}
}

\providecommand{\Schemarayon}{
\begin{circuitikz}
\draw [dash dot dot] (25.375,7.875) -- (29.75,7.875);
\draw [ rotate around={-45:(27.3125, 4.5)}] (26.5,4.75) rectangle (28.125,4.25);
\draw  (29.875,8.125) rectangle (31.5,7.625);
\node [font=\fontsize{10.2pt}{13.3pt}\selectfont, inner xsep=0.080cm, inner ysep=0.085cm, rounded corners=0.020cm] at (31.875,7.875) {Fixe};
\node [font=\fontsize{14.8pt}{19.2pt}\selectfont, inner xsep=0.080cm, inner ysep=0.085cm, rounded corners=0.020cm] at (30.75,8.5) {$UWB_1$};
\node [font=\fontsize{14.8pt}{19.2pt}\selectfont, rotate around={-45:(0,0)}, inner xsep=0.080cm, inner ysep=0.085cm, rounded corners=0.020cm] at (26.875,4) {$UWB_2$};
\draw [short] (30.875,7.125) -- (30.875,7.125);
\draw [short] (30.375,7.625) -- (30.375,7.625);
\draw [dotted] (30.75,7.5) -- (30.75,3.875);
\draw [short] (30.625,7) -- (30.625,7);
\draw [short] (30.625,6.625) -- (30.625,6.625);
\draw [densely dotted] (30.5,5.625) -- (30.5,5.625);
\draw [densely dotted] (30.625,6) -- (30.625,6);
\draw [densely dotted] (30.625,6) -- (30.625,6);
\draw [{Stealth[scale=1.5]}-{Stealth[scale=1.5]}, ] (27.625,4.75) -- (30.625,7.5);
\draw [short] (30.625,5.75) arc[start angle=-74.74, end angle=-164.74, radius=0.7126];
\node [font=\fontsize{10.2pt}{13.3pt}\selectfont, inner xsep=0.080cm, inner ysep=0.085cm, rounded corners=0.000cm] at (30,6) {$\theta$};
\draw [dash dot dot] (32.25,7.875) -- (34.625,7.875);
\node [font=\fontsize{10.2pt}{13.3pt}\selectfont, inner xsep=0.080cm, inner ysep=0.085cm, rounded corners=0.000cm] at (31,4.125) {y};
\draw [short] (29,5) -- (29,5);
\draw [short] (28.375,5.5) -- (28.375,5.5);
\draw [line width=1.2pt, short] (27.25,4.875) -- (27.75,4.375);
\draw [line width=1.2pt, short] (30.375,7.625) -- (31,7.625);
\node [font=\fontsize{14.8pt}{19.2pt}\selectfont, rotate around={41:(0,0)}, inner xsep=0.080cm, inner ysep=0.085cm, rounded corners=0.010cm] at (29,6.375) {$d$};
\draw [dash dot dot] (30.625,2.75) arc[start angle=-95.88, end angle=-130.52, radius=4.7969];
\draw [dash dot dot] (26.625,5.25) arc[start angle=-138.76, end angle=-173.32, radius=4.1447];
\end{circuitikz}
}

\providecommand{\FtFhorizontale}{
\begin{circuitikz}
\tikzstyle{every node}=[font=\fontsize{10.2pt}{13.3pt}\selectfont]
    \draw  (32.5,13.5) rectangle (34.75,13.25);
    \draw  (34.125,13.625) rectangle (34.75,13.5);
    \draw  (32.875,13.625) rectangle (33.75,13.5);
    \draw  (38.375,13.5) rectangle (40.625,13.25);
    \draw  (39.5,13.625) rectangle (40.375,13.5);
    \draw  (38.375,13.625) rectangle (39,13.5);
    \draw [{Stealth[scale=1.5]}-{Stealth[scale=1.5]}, ] (35,13.5) -- (38.125,13.5);
    \node [font=\fontsize{18.2pt}{23.7pt}\selectfont, inner xsep=0.080cm, inner ysep=0.085cm, rounded corners=0.020cm] at (36.5,13.875) {d};
    \node [font=\fontsize{14.8pt}{19.2pt}\selectfont, inner xsep=0.080cm, inner ysep=0.085cm, rounded corners=0.020cm] at (33.75,12.875) {$UWB_1$};
    \node [font=\fontsize{14.8pt}{19.2pt}\selectfont, inner xsep=0.080cm, inner ysep=0.085cm, rounded corners=0.020cm] at (39.625,12.875) {$UWB_2$};
\end{circuitikz}
}


\providecommand{\FtFVertical}{
\begin{circuitikz}
\tikzstyle{every node}=[font=\fontsize{10.2pt}{13.3pt}\selectfont]
    \node [font=\fontsize{14.8pt}{19.2pt}\selectfont, inner xsep=0.080cm, inner ysep=0.085cm, rounded corners=0.020cm] at (33.75,12.875) {$UWB_1$};
    \node [font=\fontsize{14.8pt}{19.2pt}\selectfont, inner xsep=0.080cm, inner ysep=0.085cm, rounded corners=0.020cm] at (39.625,12.875) {$UWB_2$};
    \draw [ rotate around={-90:(36.25, 11.25)}] (35.125,11.375) rectangle (37.375,11.125);
    \draw [ rotate around={-90:(36.4375, 10.9375)}] (36,11) rectangle (36.875,10.875);
    \draw [ rotate around={-90:(36.4375, 12.0625)}] (36.125,12.125) rectangle (36.75,12);
    \draw [ rotate around={-90:(40.375, 11.25)}] (39.25,11.375) rectangle (41.5,11.125);
    \draw [ rotate around={-90:(40.1875, 10.9375)}] (39.75,11) rectangle (40.625,10.875);
    \draw [ rotate around={-90:(40.1875, 12.0625)}] (39.875,12.125) rectangle (40.5,12);
    \draw [{Stealth[scale=1.5]}-{Stealth[scale=1.5]}, ] (36.75,12) -- (39.875,12);
    \node [font=\fontsize{18.2pt}{23.7pt}\selectfont, inner xsep=0.080cm, inner ysep=0.085cm, rounded corners=0.020cm] at (38.25,12.375) {d};
    \node [font=\fontsize{14.8pt}{19.2pt}\selectfont, inner xsep=0.080cm, inner ysep=0.085cm, rounded corners=0.020cm] at (41.375,11.25) {$UWB_2$};
    \node [font=\fontsize{14.8pt}{19.2pt}\selectfont, inner xsep=0.080cm, inner ysep=0.085cm, rounded corners=0.020cm] at (35.375,11.25) {$UWB_1$};
\end{circuitikz}
}

\providecommand{\FtB}{
\begin{circuitikz}
\tikzstyle{every node}=[font=\fontsize{10.2pt}{13.3pt}\selectfont]
    \draw [ rotate around={-90:(36.25, 11.25)}] (35.125,11.375) rectangle (37.375,11.125);
    \draw [ rotate around={-90:(36.4375, 10.9375)}] (36,11) rectangle (36.875,10.875);
    \draw [ rotate around={-90:(36.4375, 12.0625)}] (36.125,12.125) rectangle (36.75,12);
    \draw [{Stealth[scale=1.5]}-{Stealth[scale=1.5]}, ] (36.75,12) -- (39.875,12);
    \node [font=\fontsize{18.2pt}{23.7pt}\selectfont, inner xsep=0.080cm, inner ysep=0.085cm, rounded corners=0.020cm] at (38.25,12.375) {d};
    \node [font=\fontsize{14.8pt}{19.2pt}\selectfont, inner xsep=0.080cm, inner ysep=0.085cm, rounded corners=0.020cm] at (41.375,11.25) {$UWB_2$};
    \node [font=\fontsize{14.8pt}{19.2pt}\selectfont, inner xsep=0.080cm, inner ysep=0.085cm, rounded corners=0.020cm] at (35.375,11.25) {$UWB_1$};
    \draw [ rotate around={-90:(40.25, 11.25)}] (39.125,11.375) rectangle (41.375,11.125);
    \draw [ rotate around={-90:(40.4375, 10.9375)}] (40,11) rectangle (40.875,10.875);
    \draw [ rotate around={-90:(40.4375, 12.0625)}] (40.125,12.125) rectangle (40.75,12);
\end{circuitikz}
}
\subsubsection{Sur un rayon}

\subsection{Avant/arrière}



\vspace{0.4cm}

% --------------------------------------------------
% Bottom Decorative Rule
% --------------------------------------------------
\vfill
\hrule height 1pt

\end{document}